% !PRINTED 475
\weberpage{475}
%<<CONT
and use the formula that follows from § 34, (19):
\[
f(\omega)\,f_2\left(\frac{-3+\omega}{2}\right) = e^{-\frac{\pi i}{8}}\sqrt{2}.
\]
Accordingly, if we put
\[
f(\sqrt{-m}) = x
\]
then by (1) $x$ is the real positive root of the equation
\weq{5}{x^{24} + \gamma_2\left(\frac{-3+\sqrt{-m}}{2}\right) x^{16} - 2^8 = 0.}
For $m = 7$ we obtain by § 125, (18)
\weq{6}{f(\sqrt{-7}) = \sqrt{2},}
while in the other cases, after the values for $\gamma_2$ have been
substituted, equation (5) splits first into two factors of degree 12,
these again into two factors of degree 6, and finally these into two
of degree 3. The equations so obtained are cubic equations
for $x^8$, $x^4$, $x^2$, $x$, of which in each case we retain only the one
that has a real positive root, and which is in the end satisfied by
$f(\sqrt{-m})$ itself.

To find the first decomposition we put equation
(5) into the form:
\[
(x^{12} - a x^4)^2 - (b x^8 + c)^2 = 0,
\]
and have to determine the \gesperrt{integers} $a$, $b$, $c$ from the equations:
\[
2a + b^2 = -\gamma_2, \quad 2bc = a^2, \quad c^2 = 2^8
\]
so that $c = \pm 16$; the equation of degree 12 with
positive root therefore reads:
\[
x^{12} - b x^8 - a x^4 - 16 = 0.
\]
The numbers $a$, $b$ are easily determined from the equations above,
and in this way one finally finds the cubic equations
sought. For $m = 11$, for example, one obtains successively
\[
\begin{aligned}
x^{12} - 8x^8 + 16x^4 - 16 &= 0,\\
x^6 - 4x^2 - 4 &= 0,\\
x^3 - 2x^2 + 2x - 2 &= 0.
\end{aligned}
\]
In the five cases of § 125 one finds the following equations:
\weq{7}{x = f(\sqrt{-11}), \quad x^3 - 2x^2 + 2x - 2 = 0,}
\weq{8}{x = f(\sqrt{-19}), \quad x^3 - 2x - 2 = 0,}
\weq{9}{x = f(\sqrt{-43}), \quad x^3 - 2x^2 - 2 = 0,}
\weq{10}{x = f(\sqrt{-67}), \quad x^3 - 2x^2 - 2x - 2 = 0,}
\weq{11}{x = f(\sqrt{-163}), \quad x^3 - 6x^2 + 4x - 2 = 0.}
